\originalchapter{系综、密度矩阵与量子计数}
\label{ch:ensemble}
\source{8.333 L12--L14、L19--L23；本科气体公式仅保留高级推导所需连接。}
\originalsection{从热库到正则分布}
封闭总体系在能量壳等概率，子系统能量$E$的权重是热库态数$\Omega_B(U-E)$。展开$\log\Omega_B(U-E)=\log\Omega_B(U)-\beta E+O(E^2/C_BT^2)$，足够大的热库给$p(E)\propto\e^{-\beta E}$。下一阶定量说明有限热库误差，若子系统可交换的能量与热库热容同量级，正则近似不再充分。

经典单原子理想气$Z_N=V^N/(N!\lambda_T^{dN})$，$\lambda_T=\sqrt{2\pi\hbar^2/(m\kb T)}$。$N!$来自相同粒子不重复计数；去掉它会使熵混合悖论与不正确的非广延项出现。对不同种粒子各自除$N_\alpha!$，混合熵仍真实存在。微正则、正则、巨正则在热力学极限的局域观测可等价，需要稳定短程作用与适当可微大偏差结构；长程引力等系统不能无条件套这项等价。
\originalsection{能量与粒子数的涨落矩阵}
巨正则$\Xi=\Tr\exp[-\beta(H-\mu N)]$，以$u=-\beta,v=\beta\mu$作独立参数，则$\partial_u\log\Xi=\avg H$、$\partial_v\log\Xi=\avg N$，Hessian为$(H,N)$协方差矩阵，半正。固定$T$时$\partial_\mu\avg N=\beta\operatorname{Var}N$，结合$n^2\kappa_T=\partial n/\partial\mu$得$\operatorname{Var}N=V\kb T n^2\kappa_T$。交叉协方差解释热膨胀及不同响应之间的联系，不能一边固定$N$一边保留巨正则$N$涨落公式。
\originalsection{量子熵与变分原理}
量子统计态是正算符$\rho$、$\Tr\rho=1$，熵$S=-\kb\Tr\rho\log\rho$。canonical态$\rho_\beta=\e^{-\beta H}/Z$同时包含经典概率与本征态相干结构的可能性；热平衡在能量基对角不意味着所有非平衡态也对角。

相对熵$D(\rho\Vert\sigma)=\Tr\rho(\log\rho-\log\sigma)\ge0$可在有限维直接证明。以$\rho$本征基$\ket i$及$\sigma$本征值$s_j$写第二项，$\avg{i|\log\sigma|i}=\sum_j|\avg{i|j}|^2\log s_j\le\log\avg{i|\sigma|i}$由凹性得到；于是$D\ge\sum_ir_i\log[r_i/q_i]\ge0$，$q_i=\avg{i|\sigma|i}$。支撑不包含时相对熵为$+\infty$。

取$\sigma=\rho_\beta$，得$\Tr\rho H-TS(\rho)\ge F$，所以平衡态极小自由能泛函。再取参考$H_0$的热态$\rho_0$，得到Bogoliubov变分不等式
\begin{equation}
F\le F_0+\avg{H-H_0}_0.
\end{equation}
它对量子非对易$H,H_0$也成立；不需要错误地把$\e^{-\beta(H_0+V)}$拆成两个可交换指数。无限维需相应配分函数为trace class及期望有限。
\originalsection{分子内部自由度与量子极限}
多原子气体的平移、转动、振动在弱耦合近似下配分函数相乘。刚性线性转子$E_l=\hbar^2l(l+1)/(2I)$、简并$2l+1$，高温转动配分近$2I\kb T/\hbar^2$，低温转动态冻结；同核分子的核自旋交换可限制偶奇$l$，不能把异核结果直接套去。谐振模$Z=\e^{-\beta\hbar\omega/2}/(1-\e^{-\beta\hbar\omega})$，高温恢复每模$\kb$热容，低温指数冻结。

固体声子与黑体辐射的低温态密度积分详见既有卷；这里强调计数约束：光子与平衡声子数不守恒，$\mu=0$；原子玻色气$N$守恒，$\mu$需由数方程求。把“玻色子”当作自动$\mu=0$会使BEC数方程失去意义。
\originalsection{相同粒子的排列与交换}
两个相同粒子的Hilbert空间先取张量积，再投到对称或反对称子空间。$N$粒子的投影与配分函数为
\begin{equation}
P_\pm=\frac1{N!}\sum_{\pi\in S_N}(\pm1)^\pi U_\pi,
\qquad Z_N^\pm=\Tr\!\left[P_\pm\e^{-\beta\sum_i h_i}\right].
\end{equation}
排列分解为长度$l$循环，每个循环贡献单体迹$Z_1(l\beta)$；循环之间的组合计数给
\begin{equation}
\log\Xi_\pm=\sum_{l=1}^\infty\frac{(\pm1)^{l-1}z^l}{l}Z_1(l\beta).
\end{equation}
对自由气$Z_1(l\beta)=V/(\lambda_T^dl^{d/2})$，交换的量子virial修正由$l=2$开始。这里不需诉诸额外相互作用就得到玻色趋向聚集、费米趋向排斥的有效第二virial符号。高温$z\ll1$或$n\lambda_T^d\ll1$只保留$l=1$，恢复经典不可分辨计数。
